\section{열린 핵심, 정량적 장벽, 다음 목표}\label{LEAD:opencore}

\subsection{남은 단 하나의 명제}
\begin{conjecture}[SP1]\label{LEAD:conj-SP1}
\OPEN{} 절대상수 $C$가 있어 충분히 큰 모든 $L$과 모든 $0\le a<\Lam_L$에 대해 $a$는 $\Lam_L$의 진약수 $C\log L$개 이하의 합이다. 중복 제거 정리에 의해 반복 허용 여부는 상관없다.
\end{conjecture}
SP1은 전달정리로 \#304를 함의하고, $m=\Lam_L$에 대해 $h(m)\ll\log\log m$을 주므로 Erdős \#18의 첫 질문을 강한 형태로 해결한다. 엔트로피 하한 $H(L)\ge(1/\log2-o(1))\log L$ 때문에 SP1은 상수배를 제외하면 최적이며, 계산값 ($L\le26$에서 $H(L)/\log_2L\in[1.5,1.78]$)과 모순되지 않는다 \NUMERIC{}. 그러나 $L\le26$에서는 $\log L$과 $\log L\log\log L$을 구별할 수 없다 \HEUR{}.

\subsection{방법별 장벽}
\begin{enumerate}[leftmargin=1.6em,itemsep=3pt]
\item \textbf{탐욕(greedy).} \PROVED{} (이전) 탐욕 규칙은 $(1-o(1))\log L\log\log L$ 아래로 내려가지 못한다. 엔트로피 수준의 이상적 간격에서도 $\Theta(\log L\log\log L)$이며, 엔트로피 수준 간격 가설의 상수 $c>1$판은 거짓이다 (STRUCT S3).
\item \textbf{블록(묶음 혼합진법).} \PROVED{} (감사 대기) $\Lam_{y_j}$ 사슬의 기하적 블록은 $(\log L)^2$ 아래로 내려가지 못한다. 초지수 블록은 각 규모에서 SP1과 동치이다 (STRUCT S2).
\item \textbf{원 방법(Fourier 인증).} \PROVED{} (이전) 절댓값 경로는 창당 엔트로피의 고정 비율만큼의 지수적 절약을 요구하지만, 알려진 도구는 $L$의 다항식 크기 절약만 준다. 주호 쪽은 이번 CMM 절로 $\omega(q)=o(\sqrt L/\log L)$까지 해결되었고, $\omega(q)>L^{1/2-\eps}$는 $Y\bmod Q$의 전역 분포 문제로 남는다 (CMM).
\item \textbf{약수 간격(비공명) 방법.} \HEUR{} (리드 계산; VKGAP에서 엄밀화 중) 비공명 $\mathrm{NR}(P,T,\eta)$가 높이 $T$까지 성립하면 규모 $\ell$에서의 상대 간격은 $\log(1/\delta(\ell))\asymp\min(\log T,\ c\ell/\log L)$이고, 탐욕 사다리로 대략 $H(L)\ll L/\log T+(\log L)^2$을 준다.
 \begin{itemize}[leftmargin=1.2em]
 \item Vinogradov--Korobov: $\log T\asymp(\log L)^{3/2-\eps}$ $\Rightarrow$ $H\ll L/(\log L)^{3/2-\eps}$ (무조건적; VKGAP).
 \item 리만 가설: 명시 공식의 오차 $O(x^{1/2}\log x\log(|t|+2))$에서 $\log T\asymp\sqrt L/\log L$ $\Rightarrow$ $H\ll\sqrt L\log L$, 곧 $N(b)\ll\sqrt{\log b}\,\log\log b$로 Vose 수준에 그친다.
 \item $(\log L)^2$ (Erdős \#18)에는 $\log T\asymp M\asymp L/\log L$, 즉 RH를 넘는 지수 높이의 비공명이 필요하다. Dirichlet 동시 근사로 $T\ge(C/\eta)^M$이면 비공명은 거짓이므로 이것이 자연스러운 상한이다.
 \item 모멘트/평균값 방법의 한계: $\int_0^T|\sum_{p\in P}p^{i\tau}|^{2k}d\tau$에 Hilbert형 평균값 부등식을 쓰려면 $T$가 $k$겹 소수곱의 로그 간격의 역수보다 커야 하는데, 그 간격은 정수 간격 하한 $L^{-k}$ 수준까지 작아질 수 있다. 따라서 이 방법으로 공명 집합의 측도를 $1$ 미만으로 누르려면 $T\ge\e^{cL}$가 필요하고, 지수 높이 $\e^{cM}$의 비공명은 인증되지 않는다 (CORE에서 엄밀화 중).
 \end{itemize}
\item \textbf{매끄러운 수 경로.} \COND{문헌} $H(L)\ll L^{1-\delta}$ (작은 $\delta$). Vose를 넘으려면 $y\ge(\log x)^C$, $C<2$ 범위의 매끄러운 수 가법 정리가 필요한데, 알려진 지수는 $2$에서 멀다 (STRUCT S6).
\end{enumerate}

\subsection{다음 연구 목표}
\begin{enumerate}[leftmargin=1.6em,itemsep=2pt]
\item VK 높이를 넘는 비공명을 특정 가중 풀에 대해 얻는 방법 (평균판으로도 충분한 설계가 있는지).
\item $\omega(q)>L^{1/2-\eps}$ 법의 혼합, 곧 $Y\bmod Q$의 분포.
\item $C<2$ 범위의 매끄러운 수 가법 정리 (Vose 개선의 직접 경로).
\item 거의 모든 $a<\Lam_L$에 대한 $O(\log L)$ 덮개 (에너지 추정).
\item 이번에 증명된 초등적 핵심(중복 제거, 전달정리, cascade 조립)의 Lean~4 형식화: 저장소의 목적과 맞으나, 이 환경에서는 Mathlib 캐시가 차단되어 수행하지 못했다.
\end{enumerate}
